% Avsnitt 3.1-3.9, ur lectures/L03/appendix/a_flip_flops_and_registers.md (A.1-A.9).

\section{Från kombinatorik till sekvensnät}
\label{c3:sec:seq}

Varje krets i \chapref{c1:ch} och \ref{c2:ch} var \textbf{kombinatorisk}: dess utgång är en ren
funktion av dess nuvarande ingångar, så samma ingångar ger alltid samma utgång, och det finns
ingen föreställning om ”före” eller ”efter” inuti kretsen.

De flesta användbara system behöver mer. En räknare måste minnas sitt räknevärde. Ett register
måste hålla kvar ett värde efter att signalen som producerade det har ändrats. En lysdiod som
växlas av en knapp måste minnas om den är tänd, eftersom det inte alls ligger kodat i knappens
ögonblickliga läge.

Alla dessa behöver \textbf{minne}: en utgång som beror på kretsens egen \emph{föregående} utgång.
Ett sekvensnät är kombinatorik med en del av utgången återkopplad till ingången, och den
\textbf{återkopplingsslingan} är vad som gör ett tillståndslöst grindnät till en krets med
tillstånd. Det här kapitlet bygger den primitiv som gör kontrollerad återkoppling praktisk,
D-vippan, och sedan de två saker du bygger direkt av den: register och flankdetektorer.

% ------------------------------------------------------------------------------------------
\section{D-låset}
\label{c3:sec:latch}

Det enklaste lagringselementet är \textbf{D-låset}: två korskopplade grindar, grindade av en
\code{enable}-ingång, som lagrar en enda bit.

\begin{textdrawing}
        ┌────────────┐
   D ──►│            │──► Q
        │   D-LÅS    │
enable─►│            │──► Qn
        └────────────┘
\end{textdrawing}

\code{D} är värdet som ska lagras, \code{enable} styr om låset är \emph{öppet} (genomsläppligt)
eller \emph{stängt} (hållande), \code{Q} är det lagrade värdet och \code{Qn} är alltid dess invers.

Varje utgång återkopplas in i den andras ekvation, vilket är precis \secref{c3:sec:seq}:s
återkopplingsslinga:

\[ Q = (D' \cdot \mathit{enable} + Q_n)' \]
\[ Q_n = (D \cdot \mathit{enable} + Q)' \]

\begin{booktable}{@{}p{4em}p{12em}L@{}}
  \thead{\code{enable}} & \thead{Beteende} & \thead{Resultat} \\
  \midrule
  \code{1} & låset är \textbf{öppet} (genomsläppligt) & $Q = D$, $Q_n = D'$; utgången följer
    ingången omedelbart \\
  \code{0} & låset är \textbf{stängt} (låst) & \code{Q} och \code{Qn} håller kvar det de senast
    hade \\
\end{booktable}

\textbf{Bygg det i CircuitVerse} av en NOT-grind, två AND-grindar och två NOR-grindar (en NOR är en
OR följd av en NOT). Korskoppla återkopplingen så att den NOR som producerar \code{Q} matar den NOR
som ligger på \code{Qn}-sidan och tvärtom. Driv \code{D} och \code{enable} från
ingångsströmbrytare och betrakta \code{Q} och \code{Qn} på lysdioder.

Testa det sedan:
\begin{itemize}
  \item Med \code{enable = 1}, ändra \code{D} och bekräfta att \code{Q} följer omedelbart. Låset
    är \textbf{genomsläppligt}.
  \item Med \code{enable = 0}, bekräfta att \code{Q} är oförändrad vad \code{D} än gör. Låset är
    \textbf{låst}.
  \item Växla \code{D} några gånger medan du slår om \code{enable}, och bekräfta att låset alltid
    fryser det \code{Q} höll i det ögonblick \code{enable} gick låg. Det är vad
    \textbf{nivåkänslig} betyder: genomsläpplig hela den tid \code{enable} är hög, inte bara i ett
    enda ögonblick.
\end{itemize}

Den sista punkten är också låsets grundläggande svaghet. Eftersom det är genomsläppligt hela den
tid \code{enable = 1} går varje glitch på \code{D} under det fönstret rakt igenom till \code{Q}. I
ett stort synkront system är det långt svårare att garantera stabila data under en hel klockfas än
att garantera dem kring ett enda ögonblick. Det är problemet D-vippan löser.

% ------------------------------------------------------------------------------------------
\section{D-vippan}
\label{c3:sec:dff}

En \textbf{D-vippa} lagrar en bit på samma sätt som ett lås, men är \textbf{flanktriggad} i stället
för nivåkänslig: den samplar \code{D} bara i det \emph{ögonblick} klockan skiftar, inte under hela
den tid klockan är hög eller låg.

Mata båda med samma \code{D} och samma klocka, så syns skillnaden direkt:

\lecfig[0.78]{lectures/L03/appendix/images/latch_vs_flipflop.png}{Ett lås följer \code{D} under
varje fönster där klockan är hög, medan en vippa ändrar sig bara vid de stigande
flankerna.}{c3:fig:latchvsff}

Låset ändrar sig tre gånger: två gånger för att \code{D} rörde sig medan klockan var hög, och en
gång för att klockan gick hög mot ett \code{D} som redan hade ändrats. Vippan ändrar sig en gång,
eftersom bara en av de fyra stigande flankerna fann \code{D} skilt från det värde \code{Q} redan
höll. Båda beter sig korrekt; de svarar på olika frågor.

\begin{textdrawing}
        ┌──────────────────┐
   D ──►│                  │──► Q
clock ─►│     D-VIPPA      │
enable─►│ (stigande flank) │──► Qn
reset_n►│                  │
        └──────────────────┘
\end{textdrawing}

\code{clock} är tidsreferensen (\secref{c3:sec:clock}). \code{enable} grindar om vippan tar emot
ett nytt värde \emph{vid den triggande flanken}. \code{reset_n} är en asynkron, aktiv låg reset:
närhelst den är \code{0} tvingas \code{Q} omedelbart till \code{0}, oavsett klockan.

\begin{booktable}{@{}p{15em}L@{}}
  \thead{Villkor} & \thead{Resultat} \\
  \midrule
  \code{reset_n = 0} & $Q = 0$, $Q_n = 1$, omedelbart (asynkront, ignorerar klockan helt) \\
  \code{reset_n = 1}, stigande flank på \code{clock}, \code{enable = 1} & $Q = D$, $Q_n = D'$,
    samplat i det ögonblicket \\
  \code{reset_n = 1}, stigande flank på \code{clock}, \code{enable = 0} & \code{Q} håller kvar
    sitt föregående värde \\
  \code{reset_n = 1}, ingen stigande flank på \code{clock} & \code{Q} håller kvar sitt föregående
    värde, vad \code{D} än gör \\
\end{booktable}

Den sista raden är värd att se snarare än att läsa:

\lecfig[0.72]{lectures/L03/appendix/images/dff_timing.png}{\code{D} går hög mitt i en cykel och
fångas vid nästa stigande flank, dippar sedan och återhämtar sig helt mellan två flanker utan att
\code{Q} någonsin rör sig.}{c3:fig:dfftiming}

\code{D} dippar låg och kommer tillbaka inom en enda klockperiod, och \code{Q} rör sig aldrig,
eftersom ingen flank inträffade medan den var låg. Ett lås i samma läge hade följt med ner och upp
igen.

Invärtes är en flanktriggad vippa två D-lås i serie, ett \textbf{master-slave}-par. Det ena är
genomsläppligt medan klockan är låg och fångar \code{D}; det andra är genomsläppligt medan klockan
är hög och för det fångade värdet vidare till \code{Q}. Nettoeffekten är att \code{Q} ändrar sig
bara i det ögonblick klockan går från låg till hög, aldrig medan den ligger kvar på någondera
nivån. Det yttre beteendet är vad VHDL:s \code{rising_edge()} modellerar
(\secref{c3:sec:template}).

\textbf{Bygg den i CircuitVerse} genom att bygga vidare på ditt lås från \secref{c3:sec:latch}:
\begin{itemize}
  \item Duplicera låset så att två sitter i serie: \emph{mastern} matar \emph{slaven}, vars
    \code{Q} är vippans utgång.
  \item Lägg till \code{clock}, \code{reset_n}, \code{D} och \code{enable} som ingångar.
  \item Driv de två låsens genomsläpplighet från klockan och dess invers. För att hålla de två
    betydelserna isär, kalla den interna styrsignalen per lås för \code{gate}: masterns
    \code{gate} är \code{clock'} och slavens är \code{clock}. Vippans yttre \code{enable} är en
    annan signal och kopplas separat, nedan.
  \item Koppla det yttre \code{enable} som en 2-till-1-mux (\secref{c2:sec:mux}) på masterns
    dataingång: när \code{enable = 1} släpper muxen igenom \code{D}, och när \code{enable = 0}
    släpper den tillbaka in vippans eget \code{Q}, så att flanken fångar om det värde som redan är
    lagrat. Grinda \emph{datat}, aldrig klockan.
  \item Lägg till resetten i \textbf{båda} låsen. Att tvinga \code{Q} låg kräver två ändringar i
    varje lås, inte en:
  \begin{itemize}
    \item mata in ett aktivt högt \code{reset} (\code{reset_n} genom en NOT) i den NOR som
      producerar \code{Q}, vilket tvingar \code{Q} till \code{0}.
    \item AND:a in \code{reset_n} i termen \code{D AND gate} i samma lås, vilket frigör \code{Qn}
      att gå hög.
  \end{itemize}
  \item Sätt klockperioden till något långsamt nog att titta på, t.ex. \code{1000 ms}.
\end{itemize}

Med resetten inkopplad blir ekvationerna för varje lås:

\[ Q = (D' \cdot \mathit{gate} + Q_n + \mathit{reset})' \]
\[ Q_n = (D \cdot \mathit{gate} \cdot \mathit{reset\_n} + Q)' \]

En frestande genväg är att låta \code{Qn}-ekvationen vara och i stället bryta slingan, genom att
AND:a \code{reset_n} på den \code{Qn}-ledning som matar \code{Q}:s NOR. Det räcker inte. Med
\code{gate = 1} och \code{D = 1} lämnar det \code{Qn} på \code{0}, vilket är det som håller
\code{Q} på \code{1}, så resetten gör i tysthet ingenting i precis det fall du är mest benägen att
testa den i. Båda ändringarna, i båda låsen, eller ingendera.

Bekräfta att \code{Q} uppdateras bara i det ögonblick klockan går hög, och ligger stilla mellan
flankerna även om \code{D} ändrar sig, och att en låg \code{reset_n} tvingar \code{Q} till \code{0}
omedelbart, medan ett släpp av den inte ändrar något förrän vid nästa stigande flank.

% ------------------------------------------------------------------------------------------
\section{Klockning}
\label{c3:sec:clock}

\textbf{Klockan} är en periodisk fyrkantvåg, som varje vippa i en synkron krets läser från samma
källa. Två tal beskriver den: \textbf{perioden} $T$, tiden för en hel cykel, och \textbf{frekvensen}
$f = 1/T$ i Hz.

Varje period har exakt två \textbf{flanker}: en \textbf{stigande flank} från \code{0} till
\code{1}, och en \textbf{fallande flank} från \code{1} till \code{0}. En synkron design väljer en
av dem, nästan alltid den stigande flanken, och uppdaterar varje vippa vid den flanken och bara
den. Att använda en enda flank överallt är vad som gör tajmingen i en stor krets analyserbar: varje
vippa ändrar sig vid samma väldefinierade ögonblick, så du kan resonera en klockcykel i taget i
stället för att följa varje signal kontinuerligt.

Simulerade och verkliga klockor går i vitt skilda tidsskalor, och logiken är identisk i båda
fallen:
\begin{itemize}
  \item I CircuitVerse, sätt en period du hinner följa, t.ex. \code{1000 ms}.
  \item På DE0-CV är systemklockan fasta \code{50 MHz}, en period på \code{20 ns}. Ingen uppfattar
    enskilda flanker i den takten, så du resonerar i termer av ”vid nästa stigande flank”, aldrig i
    verklig tid.
\end{itemize}

% ------------------------------------------------------------------------------------------
\section{Register: vippor parallellt}
\label{c3:sec:register}

Ett \textbf{register} lagrar mer än en bit: N D-vippor sida vid sida, en per bit, som alla delar
samma \code{clock}, \code{reset_n} och \code{enable}.

\begin{textdrawing}
D(0) ──►[D-FF]──► Q(0)
D(1) ──►[D-FF]──► Q(1)
D(2) ──►[D-FF]──► Q(2)
   ⋮        ⋮
D(N-1)──►[D-FF]──► Q(N-1)

  (clock, reset_n, enable delas av varje vippa ovan)
\end{textdrawing}

Det är hela idén: ett 8-bitars register är åtta D-vippor som uppdateras vid samma klockflank.

I VHDL instansierar du sällan N vippor explicit. Deklarera en \code{std_logic_vector} och tilldela
hela den inuti en enda synkron process (\secref{c3:sec:template}): varje bit syntetiseras till en
egen vippa, och alla delar processens klocka och reset.

\begin{vhdlcode}
signal q: std_logic_vector(3 downto 0);
...
process (clock, reset_n) is
begin
    if (reset_n = '0') then
        q <= "0000";
    elsif (rising_edge(clock)) then
        if (enable = '1') then
            q <= d; -- d is a std_logic_vector(3 downto 0) input
        end if;
    end if;
end process;
\end{vhdlcode}

Fyra bitar, en process, fyra vippor i hårdvara: i princip inte annorlunda än fyra
\code{d_flip_flop}-instanser kopplade parallellt.

% ------------------------------------------------------------------------------------------
\section{Den synkrona processmallen i VHDL}
\label{c3:sec:template}

Varje synkront element i det här kapitlet, vippan, registret och flankdetektorn, använder samma
form. Utan reset är en enda D-vippa:

\begin{vhdlcode}
process (clock) is
begin
    if (rising_edge(clock)) then
        q <= d;
    end if;
end process;
\end{vhdlcode}

\code{q} tar \code{d}:s värde vid varje stigande flank och håller kvar det resten av tiden. Den
enda raden inuti \code{if}-satsen räcker för att syntesen ska härleda en riktig D-vippa.

\subsection*{Regler}
\begin{itemize}
  \item \textbf{\code{clock} är den enda signalen i känslighetslistan}, om det inte finns en
    asynkron reset. Processen ska köra bara när klockan ändras; den har ingen anledning att reagera
    på något annat.
  \item En process reagerar på \emph{varje} ändring hos en signal i känslighetslistan, inte bara
    stigande, så \code{rising_edge()} filtrerar specifikt ut den stigande flanken. All logik som
    uppdateras synkront ligger inuti det \code{if}:et.
  \item Varje signal som tilldelas inuti en sådan process syntetiseras till en vippa, en per bit.
\end{itemize}

\subsection*{Med en asynkron reset}
En asynkron reset måste också stå i känslighetslistan, eftersom den kan aktiveras oberoende av
klockan, och den har företräde när den är aktiv.

\begin{vhdlcode}
process (clock, reset_n) is
begin
    if (reset_n = '0') then
        q <= '0';
    elsif (rising_edge(clock)) then
        q <= d;
    end if;
end process;
\end{vhdlcode}

Läs den i prioritetsordning: kontrollera resetten först, och leta efter en stigande flank bara om
den är inaktiv. Om ingetdera gäller (en fallande flank, eller att \code{reset_n} går från \code{0}
tillbaka till \code{1}, vilket också triggar om processen) tilldelas ingenting och \code{q} behåller
sitt föregående värde.

\subsection*{När tilldelningen faktiskt får effekt}
En regel ligger under allt ovanstående, och det är det enda stället där \code{<=} beter sig olikt
tilldelning i alla språk du har skrivit mjukvara i.

En signaltilldelning får inte effekt när dess rad körs. Den \textbf{schemalägger} ett värde, som
tillämpas först när processen avslutat det här varvet och suspenderar. Fram till dess ger varje
läsning av den signalen det värde den höll \emph{innan varvet började}, hur många tilldelningar som
än redan har körts.

Det låter som en teknikalitet. Det är skälet till att mallen ovan är en vippa och inte en ledning:

\begin{vhdlcode}
process (clock) is
begin
    if (rising_edge(clock)) then
        b <= a;
        c <= b;
    end if;
end process;
\end{vhdlcode}

Läst som mjukvara slutar både \code{b} och \code{c} med att hålla \code{a}. I hårdvara gör de inte
det. Vid varje flank schemaläggs \code{b} att ta \code{a}:s värde, och \code{c} schemaläggs att ta
det värde \code{b} höll \emph{på väg in i den flanken}, inte det som just schemalagts på den.
Resultatet är två vippor i en kedja: \code{a} når \code{b} efter en flank och \code{c} efter två.
Byt plats på raderna och ingenting ändras, eftersom ingen av tilldelningarna kan observera den
andras resultat.

Två konsekvenser värda att ta med sig:
\begin{itemize}
  \item \textbf{Ordningen spelar ingen roll} mellan tilldelningar till \emph{olika} signaler i
    samma klockade process. Alla läser de gamla värdena och alla får effekt samtidigt. Det är
    motsatsen till den sekventiella läsning som ordet ”process” inbjuder till, och det är därför en
    kedja av vippor kan skrivas i vilken ordning som helst.
  \item \textbf{Tilldela samma signal två gånger i ett varv och bara den sista tilldelningen
    överlever.} De tidigare kastas utan att någonsin nå signalen.
\end{itemize}

\Secref{c3:sec:edge}:s flankdetektor är precis tvåelementskedjan ovan, och fungerar av precis det
här skälet. \Chapref{c5:ch} återvänder till regeln från andra hållet, med en \code{variable}, som
uppdateras omedelbart och därför beter sig så som den mjukvarumässiga läsningen väntar sig.

% ------------------------------------------------------------------------------------------
\section{Flankdetektering}
\label{c3:sec:edge}

Vippans andra vardagsuppgift är att detektera det \emph{ögonblick} en signal ändras, snarare än att
läsa av dess nivå: lagra signalens föregående värde i en vippa och jämför det sedan med det
nuvarande värdet med en grind.

\begin{itemize}
  \item \textbf{Detektering av stigande flank} (\code{0} till \code{1}):
    $\mathit{edge} = \mathit{current} \cdot \mathit{previous}'$
  \item \textbf{Detektering av fallande flank} (\code{1} till \code{0}):
    $\mathit{edge} = \mathit{current}' \cdot \mathit{previous}$
\end{itemize}

\begin{textdrawing}
signal ──┬───────────────────────────► AND ──► edge
         │                              ▲
         └──►[D-FF]── previous ──(NOT)──┘
                ▲
           clock, reset_n
\end{textdrawing}

\code{edge} är hög i exakt en klockcykel, den omedelbart efter övergången, eftersom
\code{previous} kommer ikapp först vid \emph{nästa} flank. Det är mekanismen bakom ”gör något en
gång per knapptryckning” i stället för ”gör något varje klockcykel knappen hålls nere”, vilket är
precis vad \secref{c3:sec:ledtoggle} löser.

% ------------------------------------------------------------------------------------------
\section{Genomarbetat exempel: en D-vippa i VHDL}
\label{c3:sec:dffvhdl}

\file{d_flip_flop.vhd} implementerar \secref{c3:sec:dff} och \ref{c3:sec:template} tillsammans: ett
enable som grindar om en stigande flank fångar \code{d}, och en asynkron aktiv låg reset.

\vhdlfile{lectures/L03/d_flip_flop/d_flip_flop.vhd}

Dess enda process är varianten med reset plus enable av \secref{c3:sec:template}:s mall: resetten
har företräde, sedan fångar en stigande flank \code{d} in i en intern signal \code{q_s} om
\code{enable = '1'}. \code{q_s} finns eftersom en \emph{utgångsport} inte kan läsas tillbaka inuti
samma arkitektur; \code{q} och \code{q_n} drivs båda kombinatoriskt från den, utanför processen.

För att prova den på hårdvara, tilldela \code{clock} till systemklockan på \code{50 MHz},
\code{reset_n} till en tryckknapp, \code{d} och \code{enable} till två skjutströmbrytare, och
\code{q}/\code{q_n} till två lysdioder. Slå sedan på \code{enable} och ändra \code{d}, så följer
lysdioden med; slå av \code{enable} och ändra \code{d}, så gör den inte det.

Det finns ingen klockflank du kan uppfatta vid \code{50 MHz}, så vad du egentligen bekräftar är att
logiken fortfarande håller när du inte längre kan se de enskilda flankerna. Vippan ser ut att svara
på din strömbrytare; i själva verket svarade den på en av femtio miljoner klockflanker den
sekunden, och strömbrytaren avgjorde bara vilket värde den flanken fångade.

% ------------------------------------------------------------------------------------------
\section{Genomarbetat exempel: flankdetekterad lysdiodsväxling}
\label{c3:sec:ledtoggle}

\file{led_toggle.vhd} kombinerar \secref{c3:sec:register}, \ref{c3:sec:template} och
\ref{c3:sec:edge} till ett litet system: två oberoende tryckknappar, som var och en växlar sin egen
lysdiod exakt en gång per tryck.

\vhdlfile{lectures/L03/led_toggle/led_toggle.vhd}

Tre 2-bitars signaler, två hållna i klockade processer och en beräknad mellan dem:
\begin{itemize}
  \item \code{button_prev_n} är ett 2-bitars register (\secref{c3:sec:register}) som håller förra
    cykelns \code{button_n}, resatt till \code{"11"} eftersom knapparna är aktiva låga.
  \item \code{button_edge <= (not button_n) and button_prev_n} är \secref{c3:sec:edge}:s detektor
    för fallande flank tillämpad på båda bitarna samtidigt, hög i en cykel exakt när en knapp går
    från släppt till nedtryckt.
  \item \code{led_s} är ett andra 2-bitars register som, i stället för att fånga ett nytt värde
    varje cykel, \textbf{växlar} bit \code{i} bara när \code{button_edge(i) = '1'}. Utan den
    grindningen skulle det växla varje klockcykel i stället för en gång per tryck, vilket är hela
    poängen med att beräkna \code{button_edge}.
\end{itemize}

\paragraph{På kortet}
Skapa Quartus Prime Lite-projektet mot DE0-CV:s \code{5CEBA4F23C7N} precis som i \chapref{c1:ch},
tilldela sedan \code{clock} till oscillatorn på \code{50 MHz}, \code{reset_n} och
\code{button_n(1 downto 0)} till tre tryckknappar, och \code{led(1 downto 0)} till två lysdioder.
Varje tryck ska växla exakt en lysdiod, en gång, hur länge du än håller den nere.

\begin{warning}
\textbf{En reservation, medvetet lämnad öppen.} \code{button_n} är en rå signal rakt från en fysisk
tryckknappspinne. I simulering är det en ren, ögonblicklig övergång; på riktig hårdvara är den
varken ögonblicklig (kontakter studsar) eller synkroniserad till \code{clock}. Den här
flankdetektorn är logiskt korrekt men \emph{osäker} att koppla direkt till en riktig knapp på en
FPGA utan mer arbete först, och det arbetet är precis vad \chapref{c4:ch} tar upp.
\end{warning}
